\section{Jahresabschlussbericht}

%___________________________________________________________________
% \begin{frame}{Bilanz}

    
    % Prognosebericht:
    % Makr für Mikro- und nanotechnologie ist gut und stabil.
    % EU möchte mehr souverinität, EU-ChipsAct
    % Stabile nachfrage der industrie and Tehcnollgien des IHPs

    % risikobericht:
    % ukrainekrieg schwer vorhersehbar. schwankende energiekosren sind das größte kostenrisiko.
    % fachkräftemangel, daher seit 2024 marketing argentur
    % drittmittel zwingend erfoderlich 

% \end{frame}

%___________________________________________________________________
\begin{frame}{Bilanz}

\begin{table}
\centering
\begin{tabular}{l S[table-format=3.1] l S[table-format=3.1]}
    \toprule
    \multicolumn{2}{c}{Aktiva in \unit{\mega\euro}} &
    \multicolumn{2}{c}{Passiva in \unit{\mega\euro}} \\
    \cmidrule(lr){1-2} \cmidrule(lr){3-4}

    Anlagevermögen & 75,4 &
    Eigenkapital & 17,0 \\
        
    \subtopic{Technische Anlagen und Maschinen} & \subnumtopic{33,3} &
    \subtopic{Jahresüberschus} & \subnumtopic{5,9} \\

    Umlaufvermögen & 46,3 &
     Sonderposten für Zuschüsse\footnote{%
        Öffentliche Fördermittel für Investitionen, die über die Nutzungsdauer der Anlagen aufgelöst werden.
    } & 78,7 \\
    
    \subtopic{Forderungen an Zuwendungsgeber} & \subnumtopic{11,0} &
    Rückstellungen & 1,6 \\

    Rechnungsabgrenzungsposten & 0,9 &
    Verbindlichkeiten & 25,4 \\
    
    \bottomrule
\end{tabular}
\caption{Bilanz des IHP zum 31.12.2024 Bilanzsumme: \qty{122,6}{\mega\euro}.}
\end{table}

\end{frame}
